\section{Step 61 Preflight two independent windows}\label{h263-block-s01}
Validate the current window using extent $7s+8$, then the prediction window
using its own stride for INTER. Check all controls and reserved fields before
reading either window. For INTRA, require zero prediction fields and omit that
window entirely. The last active byte is determined by the last row, not by a
rounded allocation size.

\paragraph{Observe} A current stride of nine and prediction stride of seventeen
are legal together. Neither row's padding is part of the input.
\Needspace{8\baselineskip}
\paragraph{Check} Which prediction stride applies during INTER admission?
\begin{itemize}
\item[A.] The independent prediction stride in the request.
\item[B.] Always the current stride.
\item[C.] Always the full QCIF width.
\end{itemize}

\section{Step 62 Gather signed transform input}\label{h263-block-s02}
Read current bytes as unsigned values. For INTER, read matching prediction bytes
and subtract prediction from current in a signed general register. For INTRA,
use zero prediction, leaving absolute current values. Store tight signed halfwords
in a private transform-input array and retain tight halfword prediction values
for the later reconstruction addition.

The subtraction range $[-255,255]$ fits the forward transform's admitted domain.
INTRA values 0 through 255 also fit. Avoid an unsigned-byte subtraction that
would replace a negative residual with its wrapped representation.

\paragraph{Observe} Current five predicted from ten produces signed residual
minus five; it remains negative when loaded as a halfword.
\Needspace{8\baselineskip}
\paragraph{Check} Why use zero prediction for INTRA input here?
\begin{itemize}
\item[A.] INTRA stores only motion vectors.
\item[B.] Its transform describes absolute samples in this baseline path.
\item[C.] It changes the quantizer parameter to zero.
\end{itemize}

\section{Step 63 Call the forward transform privately}\label{h263-block-s03}
Build the earlier transform request in aligned private storage, using input
stride sixteen, direction zero and the selected arithmetic mode. Point output
to a separate private coefficient block. Save the return address and any values
needed across the call according to the register contract.

Inspect the child status before using its output. On failure, restore saved
state and return parent status two, retaining every byte of caller output.
Do not interpret a nonzero child status as a coefficient count.

\paragraph{Observe} The constant residual two produces DC sixteen in the named
trace. That coefficient block is still private and is not the final result.
\Needspace{8\baselineskip}
\paragraph{Check} What direction belongs in this first transform request?
\begin{itemize}
\item[A.] Inverse.
\item[B.] Picture playback.
\item[C.] Forward.
\end{itemize}

\section{Step 64 Quantize and retain the ordinary count}\label{h263-block-s04}
Build the quantizer request from the private forward block, parameter, block type
and mode. Its output supplies the selected-level and reconstructed-coefficient
arrays. Save the returned ordinary nonzero count in private storage before
calling the inverse, which uses the same return registers for its own result.

The count excludes INTRA DC. It is useful for later coding decisions, but it
is not a six-block coded-pattern mask and does not yet describe the emitted bits.

\paragraph{Observe} The default INTER block selects ordinary DC level one and
returns ordinary nonzero count one.
\Needspace{8\baselineskip}
\paragraph{Check} Why save the quantizer count across the next call?
\begin{itemize}
\item[A.] The inverse call can replace the return registers.
\item[B.] It is the picture's byte stride.
\item[C.] The count changes the current source bytes in place.
\end{itemize}

\section{Step 65 Invert the reconstructed coefficient block}\label{h263-block-s05}
Point the inverse-transform request at the second quantizer output array.
Use direction one, stride sixteen and the selected mode. Its output is a third
private array of signed spatial halfwords. Check status again before reconstruction.

A pointer to the original forward block can look plausible because it has the
same size and element type. Verify the actual address in the request. Bypassing
inverse quantization would let the encoder retain information that was removed
from its transmitted representation.

\paragraph{Observe} Default reconstructed DC 23 yields inverse residual three;
using original DC sixteen would incorrectly produce residual two instead.
\Needspace{8\baselineskip}
\paragraph{Check} Which array supplies this inverse request?
\begin{itemize}
\item[A.] Original raw pixel bytes.
\item[B.] Reconstructed coefficients derived from selected levels.
\item[C.] Unquantized forward coefficients.
\end{itemize}

\section{Step 66 Add prediction before sample clipping}\label{h263-block-s06}
For scalar reconstruction, load each inverse result as a signed halfword and its
retained prediction as a nonnegative halfword. Add in a signed register, clamp
the complete sum to zero through 255, then store a byte in the private pixel
array. INTRA uses the retained zero prediction through the same arithmetic.

Keep this final sample clipping separate from the inverse's signed nine-bit
clipping and the coefficient's twelve-bit clipping. Each boundary represents a
different quantity.

\paragraph{Observe} A predicted five plus inverse residual minus ten stores zero.
Clipping the residual to zero first would incorrectly store five.
\Needspace{8\baselineskip}
\paragraph{Check} Which operation occurs first for an INTER pixel?
\begin{itemize}
\item[A.] Narrow the signed residual to an unsigned byte.
\item[B.] Clip the residual alone to the final pixel range.
\item[C.] Add signed residual to prediction in a sufficient width.
\end{itemize}

\section{Step 67 Perform bounded vector sample reconstruction}\label{h263-block-s07}
In mode one, load eight inverse and eight prediction halfwords from aligned
private arrays. Add halfword lanes, apply signed maximum with zero and signed
minimum with 255, and store exactly eight pixel bytes. A straightforward taught
version stores the clipped vector temporarily and extracts its native halfwords
for byte stores; another correct packing schedule is acceptable.

The sum bound $[-256,510]$ establishes that an ordinary halfword addition is
safe before clipping. This argument depends on the earlier inverse contract.
Test mixed signs and positions rather than only a constant vector.

\paragraph{Observe} The scalar and vector paths preserve signed inverse values
in the diagnostic array while producing clipped pixel bytes in the final array.
\Needspace{8\baselineskip}
\paragraph{Check} What must happen before narrowing a lane to a pixel byte?
\begin{itemize}
\item[A.] Clip the complete signed sum to 0 through 255.
\item[B.] Discard all negative inverse outputs before adding prediction.
\item[C.] Replace the selected levels with raw pixels.
\end{itemize}

\section{Step 68 Publish a complete block result}\label{h263-block-s08}
After every dependency and pixel operation succeeds, copy the selected levels,
reconstructed coefficients, inverse values and final pixels into the caller's
448-byte record. Return the saved ordinary count. A child can fail after writing
its own private output; none of that work authorizes an early public store.

Stop the ordinary caller at \texttt{block\_caller\_return}, before its first
instruction. The complete record and return registers are ready, but named
status/count fields still contain their old values. The default prints first
pixel 131 and exits zero after the caller resumes.

\paragraph{Observe} The default successful count is one. INTRA input 128 instead
produces DC code 255, reconstructed DC 1024, pixels 128 and ordinary count zero.
\Needspace{8\baselineskip}
\paragraph{Check} What survives a failed inverse child?
\begin{itemize}
\item[A.] A partially replaced caller record.
\item[B.] Every old caller-result byte, with parent failure status two.
\item[C.] A successful status because forward transform already completed.
\end{itemize}

\section{Step 69 Check the actual selected-level reconstruction}\label{h263-block-s09}
Run the public feedback on a source containing your block entry and completed
transform/quantizer dependencies:
\begin{verbatim}
python3 drafts/h263-reconstruction-construction/grade.py \
    my-block.asm --cli /path/to/mars-cli.jar \
    --report block-results.json
\end{verbatim}
Compare every stage, not just the first final pixel. The checker accepts the
permitted forward and inverse accuracy while requiring exact level-to-coefficient
and inverse-to-pixel consistency. It also checks wide pointers, independent
strides, input-read bounds, output guards and saved state.

The statistical inverse gate remains a separate prerequisite. An implementation
that passes one constant block can still have a normalization, transpose or bias
error elsewhere.

\paragraph{Observe} Two independently accurate precisions and another traversal
order can satisfy the same block interface without sharing every rounded coefficient.
\Needspace{8\baselineskip}
\paragraph{Check} Why reconstruct the reference inverse from returned levels?
\begin{itemize}
\item[A.] It forces every encoder to emit the host encoder's bits.
\item[B.] It avoids checking quantization entirely.
\item[C.] It checks the actual representation this block selected.
\end{itemize}

\section{Step 70 Retain reconstructed pixels for the next prediction}\label{h263-block-s10}
Use the first reconstructed block as prediction for a second call in the same
session. With constant current 132 and prior reconstruction 131, observe zero
ordinary levels and reconstructed pixels 131 at $Q=8$. Check that the second
call builds fresh requests and counts rather than reusing stale values.

The next picture-level integration must copy each completed block into its proper
Y, Cb or Cr plane and publish the new reference only after the picture succeeds.
Keeping the block result correct is necessary, but does not itself implement
picture ownership, file output, a coded macroblock or a conforming stream.

\paragraph{Observe} Reusing raw 130 instead of reconstructed 131 changes the
second residual and loses synchronization with a decoder's available reference.
\Needspace{8\baselineskip}
\paragraph{Check} Which stored pixels can become the next reference?
\begin{itemize}
\item[A.] The completed locally reconstructed picture.
\item[B.] The raw input regardless of quantization loss.
\item[C.] Any half-written pixel array after a child failure.
\end{itemize}
